\documentclass[11pt,a4paper]{article}
\usepackage[margin=1in]{geometry}
\usepackage{amsmath,amssymb}
\usepackage{booktabs,tabularx,longtable,array}
\usepackage{enumitem}
\usepackage{titlesec}
\usepackage{fancyhdr}
\usepackage[table]{xcolor}
\usepackage[hidelinks]{hyperref}
\usepackage{tcolorbox}
\tcbuselibrary{skins,breakable}

% Setup page style
\pagestyle{fancy}
\fancyhf{}
\fancyhead[L]{\textbf{Kathmandu University BDS 1st Year}}
\fancyhead[R]{\textbf{Biochemistry Solutions}}
\fancyfoot[C]{\thepage}
\renewcommand{\headrulewidth}{0.6pt}

% Section styling
\titleformat{\section}{\Large\bfseries\color{green!50!black}}{\thesection}{1em}{}[\titlerule]
\titleformat{\subsection}{\large\bfseries\color{black!85}}{\thesubsection}{1em}{}
\titleformat{\subsubsection}{\normalsize\bfseries\color{black!75}}{\thesubsubsection}{1em}{}

% Custom boxes
\newtcolorbox{qbox}[2][]{
  colback=green!5!white,
  colframe=green!60!black,
  fonttitle=\bfseries,
  title={#2},
  arc=2mm,
  breakable,
  #1
}

\newtcolorbox{ansbox}[1][]{
  colback=white,
  colframe=gray!50!black,
  fonttitle=\bfseries,
  title={Answer},
  arc=1mm,
  breakable,
  #1
}

\begin{document}

\begin{titlepage}
    \centering
    \vspace*{2cm}
    {\Huge \textbf{Kathmandu University}}\\[0.5cm]
    {\LARGE \textbf{Bachelor of Dental Surgery (BDS) 1st Year}}\\[1.5cm]
    {\Huge \textbf{BIOCHEMISTRY}}\\[0.8cm]
    {\Large \textbf{Comprehensive Exam Solutions \& Question Bank}}\\[0.5cm]
    {\large Complete Question Papers, Exam-Oriented Answers, Frequency Analysis \& Revision Cheatsheet}\\[2cm]
    \vfill
    {\large \textbf{Compiled from Past Papers: 2015--2022}}\\[0.5cm]
    {\small Kathmandu University School of Medical Sciences (KUSMS) \& Affiliated Colleges}
    \vspace*{1.5cm}
\end{titlepage}

\tableofcontents
\newpage

\section{Past Examination Questions and Solutions (2015--2022)}

\subsection{Examination: May 10, 2022 (End Semester Exam)}
\textbf{Subject:} Biochemistry \quad \textbf{Marks:} 50 (Section B: 25, Section C: 25)

\subsubsection{Section B [25 Marks]}

\begin{qbox}{Question 1 [1+3+2 = 6 Marks]}
List the components of the electron transport chain (ETC). Explain the chemiosmotic hypothesis of oxidative phosphorylation. Add a note on uncouplers.
\end{qbox}

\begin{ansbox}
\textbf{1. Components of the Electron Transport Chain (ETC):}
Located in the inner mitochondrial membrane; organized into four multi-protein complexes and two mobile electron carriers:
\begin{itemize}[leftmargin=*]
    \item \textbf{Complex I (NADH-Q Oxidoreductase / NADH Dehydrogenase):} Flavoprotein containing FMN and Fe-S clusters; accepts electrons from NADH, transfers them to Coenzyme Q, and pumps **$4\ \text{H}^+$** into intermembrane space.
    \item \textbf{Complex II (Succinate-Q Reductase / Succinate Dehydrogenase):} Contains FAD and Fe-S centers; oxidizes succinate to fumarate in TCA cycle and transfers electrons to Coenzyme Q; **does not pump protons**.
    \item \textbf{Coenzyme Q (Ubiquinone):} Lipid-soluble mobile carrier residing in the lipid bilayer; collects electrons from Complexes I and II and shuttles them to Complex III.
    \item \textbf{Complex III (Q-Cytochrome c Oxidoreductase / Cytochrome $bc_1$ Complex):} Contains Cytochromes $b_L, b_H$, Fe-S Rieske center, and Cytochrome $c_1$; transfers electrons to Cytochrome $c$ and pumps **$4\ \text{H}^+$** per electron pair.
    \item \textbf{Cytochrome c:} Water-soluble, peripheral mobile hemoprotein in the intermembrane space; shuttles single electrons from Complex III to Complex IV.
    \item \textbf{Complex IV (Cytochrome c Oxidase):} Contains Cytochromes $a, a_3$ and copper centers ($Cu_A, Cu_B$); transfers 4 electrons to molecular oxygen ($\text{O}_2$) to form $2\ \text{H}_2\text{O}$; pumps **$2\ \text{H}^+$** into intermembrane space.
\end{itemize}

\textbf{2. Chemiosmotic Hypothesis (Peter Mitchell, 1961):}
\begin{enumerate}[leftmargin=*]
    \item As high-energy electrons flow down the respiratory chain from NADH/FADH$_2$ to $\text{O}_2$, the free energy released drives active extrusion of protons ($\text{H}^+$) from the mitochondrial matrix across the inner mitochondrial membrane into the intermembrane space via Complexes I, III, and IV.
    \item The inner mitochondrial membrane is strictly impermeable to protons. Thus, proton pumping establishes a steep **electrochemical proton gradient ($\Delta \mu_{\text{H}^+}$)** across the membrane, comprising both an electrical membrane potential ($\Delta \Psi$, inside negative) and a pH gradient ($\Delta \text{pH}$, inside alkaline). This gradient is called the **Proton-Motive Force (PMF)**.
    \item Protons flow thermodynamically downhill back into the mitochondrial matrix exclusively through the proton-conducting channel ($F_0$ subunit) of **ATP Synthase (Complex V)**.
    \item The kinetic energy of proton translocation rotates the $c$-ring of $F_0$, inducing conformational transitions in the catalytic $\beta$-subunits of the $F_1$ headpiece ("binding change mechanism" of Paul Boyer), driving the phosphorylation of ADP by inorganic phosphate: $\text{ADP} + \text{P}_i \rightarrow \text{ATP}$.
\end{enumerate}

\textbf{3. Uncouplers of Oxidative Phosphorylation:}
\begin{itemize}[leftmargin=*]
    \item \textit{Mechanism:} Lipophilic weak acids that readily cross the inner mitochondrial membrane in both protonated and deprotonated states. They short-circuit the membrane by translocating protons back into the matrix, **dissipating the electrochemical proton gradient as heat without producing ATP**. Electron transport accelerates to maximal rate while ATP synthesis is abolished.
    \item \textit{Examples:}
    \begin{enumerate}
        \item \textbf{2,4-Dinitrophenol (2,4-DNP):} Classic synthetic chemical uncoupler; causes hyperthermia and metabolic collapse.
        \item \textbf{Thermogenin (UCP-1):} Physiological uncoupling protein in brown adipose tissue; generates non-shivering thermogenesis in neonates.
        \item High-dose Salicylates (Aspirin).
    \end{enumerate}
\end{itemize}
\end{ansbox}

\begin{qbox}{Question 2 [1+2+2 = 5 Marks]}
What are ketone bodies? Name them. Describe their synthesis and utilization. In which conditions is their synthesis increased?
\end{qbox}

\begin{ansbox}
\textbf{1. Definition and Names:}
\textbf{Ketone bodies} are water-soluble organic compounds synthesized by liver mitochondria from Acetyl-CoA during periods of high fatty acid oxidation. They serve as an essential alternative circulating fuel for extrahepatic tissues. The three ketone bodies are:
\begin{enumerate}[leftmargin=*]
    \item \textbf{Acetoacetate} (the primary metabolic ketone body).
    \item \textbf{$\beta$-Hydroxybutyrate} (chemically a hydroxy acid; quantitatively the predominant form in blood).
    \item \textbf{Acetone} (volatile, non-metabolizable side-product excreted via breath; causes characteristic fruity odor in ketoacidosis).
\end{enumerate}

\textbf{2. Ketogenesis (Synthesis in Liver Mitochondria):}
\begin{enumerate}[leftmargin=*]
    \item $2\ \text{Acetyl-CoA} \xrightarrow{\text{Thiolase}} \text{Acetoacetyl-CoA} + \text{CoA}$.
    \item $\text{Acetoacetyl-CoA} + \text{Acetyl-CoA} \xrightarrow{\text{\textbf{HMG-CoA Synthase (Rate-Limiting)}}} \beta\text{-Hydroxy-}\beta\text{-methylglutaryl-CoA (HMG-CoA)} + \text{CoA}$.
    \item $\text{HMG-CoA} \xrightarrow{\text{HMG-CoA Lyase}} \text{\textbf{Acetoacetate}} + \text{Acetyl-CoA}$.
    \item Fates of Acetoacetate:
    \begin{itemize}
        \item Reduced reversibly to **$\beta$-hydroxybutyrate** by $\beta$-hydroxybutyrate dehydrogenase (using NADH).
        \item Non-enzymatic spontaneous decarboxylation yields **Acetone** + $\text{CO}_2$.
    \end{itemize}
\end{enumerate}

\textbf{3. Ketolysis (Utilization by Extrahepatic Tissues):}
\begin{itemize}[leftmargin=*]
    \item Utilized by cardiac muscle, skeletal muscle, renal cortex, and the **brain during prolonged fasting/starvation** (saves muscle protein catabolism).
    \item \textit{Reactivation:} Acetoacetate is activated to Acetoacetyl-CoA by the enzyme **Thiophorase ($\beta$-ketoacyl-CoA transferase)** using Succinyl-CoA from the TCA cycle.
    \item Thiolase then cleaves Acetoacetyl-CoA into $2\ \text{Acetyl-CoA}$, which enter the TCA cycle to yield ATP.
    \item \textbf{Crucial Rule:} The **liver cannot utilize ketone bodies** because it lacks the enzyme **Thiophorase**.
\end{itemize}

\textbf{4. Conditions with Increased Ketone Body Synthesis (Ketosis):}
\begin{enumerate}[leftmargin=*]
    \item \textbf{Uncontrolled Diabetes Mellitus (Diabetic Ketoacidosis - DKA):} Insulin deficiency and glucagon excess stimulate massive adipose lipolysis, flooding the liver with free fatty acids. Oxaloacetate is depleted by excessive gluconeogenesis, diverting excess Acetyl-CoA into ketogenesis.
    \item \textbf{Prolonged Starvation / Fasting:} Carbohydrate deprivation forces tissues to rely on fatty acids and ketone bodies.
    \item High-fat, low-carbohydrate (ketogenic) diets and persistent starvation vomiting in children.
\end{enumerate}
\end{ansbox}

\begin{qbox}{Question 3 [2+3 = 5 Marks]}
Write down the irreversible steps of glycolysis with regulatory enzymes. Add a note on the Rapoport-Luebering cycle and its significance in RBCs.
\end{qbox}

\begin{ansbox}
\textbf{1. Irreversible Steps of Glycolysis:}
Glycolysis features three thermodynamically irreversible, highly exergonic reactions that serve as major regulatory checkpoints:
\begin{enumerate}[leftmargin=*]
    \item \textbf{Reaction 1 (Phosphorylation of Glucose):}
    $$\text{Glucose} + \text{ATP} \xrightarrow{\text{\textbf{Hexokinase / Glucokinase}}} \text{Glucose-6-Phosphate} + \text{ADP}$$
    \begin{itemize}
        \item *Hexokinase:* Present in all tissues; low $K_m$ (high affinity); inhibited by product Glucose-6-P.
        \item *Glucokinase:* In liver and pancreatic $\beta$-cells; high $K_m$ (active post-prandially); stimulated by insulin.
    \end{itemize}
    \item \textbf{Reaction 3 (Committed Rate-Limiting Step):}
    $$\text{Fructose-6-Phosphate} + \text{ATP} \xrightarrow{\text{\textbf{Phosphofructokinase-1 (PFK-1)}}} \text{Fructose-1,6-Bisphosphate} + \text{ADP}$$
    \begin{itemize}
        \item Allosterically activated by **Fructose-2,6-bisphosphate** (most potent) and AMP.
        \item Allosterically inhibited by high ATP and Citrate.
    \end{itemize}
    \item \textbf{Reaction 10 (Substrate-Level Phosphorylation):}
    $$\text{Phosphoenolpyruvate (PEP)} + \text{ADP} \xrightarrow{\text{\textbf{Pyruvate Kinase}}} \text{Pyruvate} + \text{ATP}$$
    \begin{itemize}
        \item Activated by Fructose-1,6-bisphosphate (feed-forward activation).
        \item Inhibited by ATP, Alanine, and glucagon-mediated phosphorylation.
    \end{itemize}
\end{enumerate}

\textbf{2. The Rapoport-Luebering Cycle in Erythrocytes (RBCs):}
\begin{itemize}[leftmargin=*]
    \item A supplementary metabolic shunt unique to mature red blood cells that bypasses the phosphoglycerate kinase step of glycolysis:
    \begin{enumerate}
        \item 1,3-Bisphosphoglycerate (1,3-BPG) is converted into **2,3-Bisphosphoglycerate (2,3-BPG)** by the bifunctional enzyme **Bisphosphoglycerate Mutase**.
        \item 2,3-BPG is subsequently hydrolyzed by **2,3-Bisphosphoglycerate Phosphatase** to 3-Phosphoglycerate, which re-enters the main glycolytic pathway.
    \end{enumerate}
    \item \textit{Energetic Cost:} Bypassing the 1,3-BPG $\rightarrow$ 3-PG step forfeits the synthesis of 1 ATP per triose (yields 0 net ATP), but is essential for respiratory physiology.
    \item \textit{Physiological Significance:}
    \begin{itemize}
        \item **2,3-BPG binds specifically to the central cavity of deoxygenated Hemoglobin (HbA)**, stabilizing the low-affinity T-state (tense state).
        \item This markedly **decreases hemoglobin's affinity for oxygen, shifting the oxygen-hemoglobin dissociation curve to the RIGHT**.
        \item Promotes efficient unloading of oxygen from oxyhemoglobin to peripheral tissues.
        \item 2,3-BPG levels increase adaptively in chronic hypoxia, high altitude, and chronic obstructive pulmonary disease.
    \end{itemize}
\end{itemize}
\end{ansbox}

\begin{qbox}{Question 4 [3 $\times$ 3 = 9 Marks]}
Write short notes on:
(a) Free radicals and antioxidant defense systems \\
(b) Classification and functions of lipoproteins \\
(c) Competitive vs. Non-competitive enzyme inhibition
\end{qbox}

\begin{ansbox}
\textbf{(a) Free Radicals & Antioxidants:}
\begin{itemize}[leftmargin=*]
    \item \textit{Free Radicals:} Chemical species possessing one or more unpaired electrons in their outer valence orbital, making them highly reactive and unstable. Include **Reactive Oxygen Species (ROS)**: Superoxide anion ($O_2^{\bullet-}$), Hydroxyl radical ($OH^{\bullet}$ - most reactive), Hydrogen peroxide ($H_2O_2$). Cause membrane lipid peroxidation, protein oxidation, and DNA strand breaks.
    \item \textit{Antioxidant Defenses:}
    \begin{enumerate}
        \item \textbf{Enzymatic Antioxidants:}
        \begin{itemize}
            \item *Superoxide Dismutase (SOD):* $2\ O_2^{\bullet-} + 2\ \text{H}^+ \rightarrow H_2O_2 + O_2$.
            \item *Catalase:* $2\ H_2O_2 \rightarrow 2\ H_2O + O_2$.
            \item *Glutathione Peroxidase (Selenium-dependent):* Neutralizes $H_2O_2$ to $H_2O$ by oxidizing reduced glutathione (GSH $\rightarrow$ GSSG).
        \end{itemize}
        \item \textbf{Non-Enzymatic Antioxidants:} **Vitamin E (Alpha-Tocopherol)** (chain-breaking lipid-soluble antioxidant in membranes), **Vitamin C (Ascorbic acid)** (water-soluble radical scavenger), Glutathione, Uric acid, and Carotenoids.
    \end{enumerate}
\end{itemize}

\textbf{(b) Classification and Functions of Lipoproteins:}
\begin{itemize}[leftmargin=*]
    \item Spherical macromolecular complexes consisting of a neutral hydrophobic non-polar lipid core (triglycerides and cholesteryl esters) enveloped by an amphipathic shell of phospholipids, free cholesterol, and apolipoproteins:
    \begin{enumerate}
        \item \textbf{Chylomicrons (Apo B-48, C-II, E):} Lowest density; transport **dietary (exogenous) triglycerides** from small intestine to peripheral tissues (adipose/muscle via Lipoprotein Lipase).
        \item \textbf{VLDL (Apo B-100, C-II, E):} Transports **endogenous triglycerides** synthesized by the liver to peripheral tissues.
        \item \textbf{LDL (Apo B-100):} Formed from VLDL via IDL; rich in cholesterol; delivers cholesterol to peripheral extrahepatic cells via LDL receptor endocytosis. **Atherogenic ("Bad Cholesterol")**.
        \item \textbf{HDL (Apo A-I, C-II, E):} Synthesized by liver and intestine; mediates **Reverse Cholesterol Transport** (scavenges excess cholesterol from peripheral tissues and foam cells and returns it to liver for biliary excretion). **Cardioprotective ("Good Cholesterol")**.
    \end{enumerate}
\end{itemize}

\textbf{(c) Competitive vs. Non-Competitive Enzyme Inhibition:}
\begin{center}
\begin{tabularx}{\textwidth}{p{3.2cm}X X}
\toprule
\textbf{Feature} & \textbf{Competitive Inhibition} & \textbf{Non-Competitive Inhibition} \\
\midrule
\textbf{Inhibitor Structure} & Structurally closely resembles substrate. & Chemically unrelated to substrate. \\
\textbf{Binding Site} & Binds directly to the \textbf{active catalytic site}. & Binds to a distinct \textbf{allosteric site}. \\
\textbf{Effect on $V_{max}$} & \textbf{Unchanged} (high substrate concentration overcomes inhibition). & \textbf{Decreased} (cannot be overcome by increasing substrate). \\
\textbf{Effect on $K_m$} & \textbf{Increased} (affinity appears decreased). & \textbf{Unchanged} (affinity is unaffected). \\
\textbf{Lineweaver-Burk Plot} & Lines intersect on the y-axis ($1/V_{max}$ unchanged; slope increases). & Lines intersect on the negative x-axis ($-1/K_m$ unchanged; y-intercept shifts upward). \\
\textbf{Clinical Examples} & Statins (inhibit HMG-CoA reductase); Methotrexate (inhibits DHFR); Allopurinol. & Cyanide (inhibits Cytochrome oxidase); Heavy metals (lead, mercury binding $-SH$ groups). \\
\bottomrule
\end{tabularx}
\end{center}
\end{ansbox}

\subsubsection{Section C [25 Marks]}

\begin{qbox}{Question 5 [6 Marks]}
Write down the steps of the TCA (Krebs) cycle. Explain why it is called the common terminal metabolic pathway and an amphibolic pathway.
\end{qbox}

\begin{ansbox}
\textbf{1. Reactions of the TCA (Citric Acid / Krebs) Cycle (Mitochondrial Matrix):}
\begin{enumerate}[leftmargin=*]
    \item \textbf{Condensation:} $\text{Acetyl-CoA (2C)} + \text{Oxaloacetate (4C)} + \text{H}_2\text{O} \xrightarrow{\text{\textbf{Citrate Synthase}}} \text{Citrate (6C)} + \text{CoA}$.
    \item \textbf{Isomerization:} $\text{Citrate} \xrightarrow{\text{Aconitase}} \text{Isocitrate (6C)}$.
    \item \textbf{First Oxidative Decarboxylation:} $\text{Isocitrate} + \text{NAD}^+ \xrightarrow{\text{\textbf{Isocitrate Dehydrogenase}}} \alpha\text{-Ketoglutarate (5C)} + \text{NADH} + \text{H}^+ + \text{CO}_2$.
    \item \textbf{Second Oxidative Decarboxylation:} $\alpha\text{-Ketoglutarate} + \text{NAD}^+ + \text{CoA} \xrightarrow{\text{\textbf{$\alpha$-Ketoglutarate Dehydrogenase Complex}}} \text{Succinyl-CoA (4C)} + \text{NADH} + \text{H}^+ + \text{CO}_2$.
    \item \textbf{Substrate-Level Phosphorylation:} $\text{Succinyl-CoA} + \text{GDP} + \text{P}_i \xrightarrow{\text{Succinate Thiokinase}} \text{Succinate (4C)} + \text{GTP} + \text{CoA}$.
    \item \textbf{Dehydrogenation:} $\text{Succinate} + \text{FAD} \xrightarrow{\text{Succinate Dehydrogenase (Complex II)}} \text{Fumarate (4C)} + \text{FADH}_2$.
    \item \textbf{Hydration:} $\text{Fumarate} + \text{H}_2\text{O} \xrightarrow{\text{Fumarase}} \text{L-Malate (4C)}$.
    \item \textbf{Regeneration of Oxaloacetate:} $\text{L-Malate} + \text{NAD}^+ \xrightarrow{\text{Malate Dehydrogenase}} \text{Oxaloacetate (4C)} + \text{NADH} + \text{H}^+$.
\end{enumerate}
\textit{Net Yield per Acetyl-CoA:} $3\ \text{NADH} + 1\ \text{FADH}_2 + 1\ \text{GTP} + 2\ \text{CO}_2 \rightarrow \mathbf{10\ \text{ATP}}$ (via oxidative phosphorylation).

\textbf{2. Why TCA is the "Common Terminal Metabolic Pathway":}
Carbohydrates (via pyruvate), Lipids (via fatty acid $\beta$-oxidation), and Proteins (glucogenic/ketogenic amino acid transamination) all converge biochemically to form **Acetyl-CoA** (or intermediate metabolites of the cycle). The TCA cycle is the final shared pathway where carbon skeletons from all classes of food are completely oxidized to $\text{CO}_2$ and $\text{H}_2\text{O}$.

\textbf{3. Why TCA is an "Amphibolic Pathway":}
Serves dual functions, acting simultaneously in **catabolic** and **anabolic** processes:
\begin{itemize}[leftmargin=*]
    \item \textit{Catabolic Role:} Oxidative breakdown of Acetyl-CoA to generate reducing equivalents (NADH, FADH$_2$) for energy production.
    \item \textit{Anabolic Role (Precursor Reservoir):}
    \begin{itemize}
        \item **Succinyl-CoA:** Essential substrate for **Heme biosynthesis**.
        \item **Oxaloacetate:** Transaminated to Aspartate (for purine/pyrimidine synthesis) or utilized for **Gluconeogenesis**.
        \item **$\alpha$-Ketoglutarate:** Transaminated to Glutamate, Glutamine, Proline, and GABA.
        \item **Citrate:** Transported out of mitochondria into cytosol to provide Acetyl-CoA for **Fatty Acid and Cholesterol synthesis**.
    \end{itemize}
\end{itemize}
\end{ansbox}

\begin{qbox}{Question 6 [6 Marks]}
Discuss diagrammatically the absorption of non-heme iron, its transport, and storage. How does hepcidin regulate iron absorption?
\end{qbox}

\begin{ansbox}
\textbf{1. Absorption of Non-Heme Iron (Duodenum and Upper Jejunum):}
\begin{itemize}[leftmargin=*]
    \item Dietary non-heme iron exists predominantly as insoluble ferric iron ($\text{Fe}^{3+}$).
    \item At the apical enterocyte brush border, **Duodenal Cytochrome b (Dcytb)** (a ferrireductase stimulated by Vitamin C/gastric acid) reduces insoluble $\text{Fe}^{3+}$ to absorbable ferrous iron ($\text{Fe}^{2+}$).
    \item $\text{Fe}^{2+}$ is transported into the enterocyte cytoplasm via the **Divalent Metal Transporter-1 (DMT-1)** ($H^+$-coupled symporter).
\end{itemize}

\textbf{2. Basolateral Export and Blood Transport:}
\begin{itemize}[leftmargin=*]
    \item Intracellular $\text{Fe}^{2+}$ is exported across the basolateral membrane into the bloodstream by the transmembrane iron efflux channel **Ferroportin-1**.
    \item On the capillary side, the ferroxidase **Hephaestin** (and circulating ceruloplasmin) oxidizes $\text{Fe}^{2+}$ back to $\text{Fe}^{3+}$.
    \item In plasma, $\text{Fe}^{3+}$ binds tightly to the transport glycoprotein **Transferrin** (each transferrin molecule binds $2\ \text{Fe}^{3+}$ ions). Normal transferrin saturation is 20--45\%.
\end{itemize}

\textbf{3. Intracellular Iron Storage:}
\begin{itemize}[leftmargin=*]
    \item Excess iron inside cells (hepatocytes, reticuloendothelial macrophages of spleen/bone marrow, enterocytes) is stored in the protein shell of **Ferritin** (apoferritin + ferric phosphate core, storing up to 4,500 iron atoms).
    \item When iron stores are excessively overloaded, ferritin aggregates into insoluble, visible golden-brown **Hemosiderin** granules.
\end{itemize}

\textbf{4. Regulation by Hepcidin (Master Iron Regulator):}
\begin{itemize}[leftmargin=*]
    \item \textit{Origin:} Small 25-amino acid peptide hormone synthesized and secreted by the **liver**.
    \item \textit{Mechanism of Action:} Hepcidin binds directly to the cellular iron exporter **Ferroportin-1** on the basolateral surface of duodenal enterocytes and reticuloendothelial macrophages.
    \item This induces internalization, ubiquitination, and lysosomal degradation of ferroportin.
    \item \textit{Consequences:}
    \begin{itemize}
        \item Enterocytes cannot export absorbed iron into blood; trapped iron is lost when enterocytes desquamate (reduces dietary iron absorption).
        \item Macrophages cannot release recycled iron from senescent RBCs into plasma.
    \end{itemize}
    \item \textit{Regulation of Hepcidin Secretion:} Upregulated by high iron stores and inflammatory cytokines (**IL-6** in chronic infections/autoimmunity, leading to **Anemia of Chronic Disease**); downregulated by iron deficiency, hypoxia, and active erythropoiesis.
\end{itemize}
\end{ansbox}

\begin{qbox}{Question 7 [5 Marks]}
Define neurotransmitter. Write down the criteria for considering a substance as a neurotransmitter and classify them. Outline catecholamine synthesis.
\end{qbox}

\begin{ansbox}
\textbf{1. Definition:}
A \textbf{neurotransmitter} is an endogenous chemical messenger synthesized within a neuron, packaged into synaptic vesicles, and released into the synaptic cleft upon depolarization, binding to specific postsynaptic receptors to alter the electrical or biochemical activity of the target cell.

\textbf{2. Classical Criteria to Qualify as a Neurotransmitter:}
\begin{enumerate}[leftmargin=*]
    \item Must be synthesized within the presynaptic neuron.
    \item Must be present in presynaptic terminals in sufficient concentrations to exert a physiological effect.
    \item Must be released upon presynaptic depolarization in a $Ca^{2+}$-dependent exocytotic mechanism.
    \item Exogenous application to the postsynaptic membrane must mimic the exact response produced by endogenous nerve stimulation.
    \item Specific mechanisms for rapid inactivation or termination must exist (enzymatic breakdown or cellular reuptake).
\end{enumerate}

\textbf{3. Classification of Neurotransmitters:}
\begin{enumerate}[leftmargin=*]
    \item \textbf{Small-Molecule Neurotransmitters:}
    \begin{itemize}
        \item *Acetylcholine (ACh):* Neuromuscular junction, parasympathetic nerves.
        \item *Biogenic Amines (Catecholamines \& Indolamines):* Dopamine, Norepinephrine, Epinephrine, Serotonin (5-HT), Histamine.
        \item *Amino Acids:*
        \begin{itemize}
            \item Excitatory: **Glutamate** (primary in CNS), Aspartate.
            \item Inhibitory: **GABA** (Gamma-Aminobutyric Acid - primary in brain), **Glycine** (primary in spinal cord).
        \end{itemize}
        \item *Purines:* ATP, Adenosine.
    \end{itemize}
    \item \textbf{Neuropeptides (Large-Molecule):} Substance P (pain), Endorphins/Enkephalins, Neuropeptide Y, VIP.
    \item \textbf{Gas Neurotransmitters:} Nitric Oxide (NO), Carbon Monoxide (CO).
\end{enumerate}

\textbf{4. Pathway of Catecholamine Biosynthesis:}
Occurs in the adrenal medulla and adrenergic neurons from the amino acid **L-Tyrosine**:
\begin{enumerate}[leftmargin=*]
    \item $\text{L-Tyrosine} \xrightarrow{\text{\textbf{Tyrosine Hydroxylase (Rate-Limiting)}}} \text{\textbf{L-DOPA}}$ (requires Tetrahydrobiopterin - $BH_4$).
    \item $\text{L-DOPA} \xrightarrow{\text{DOPA Decarboxylase}} \text{\textbf{Dopamine}} + \text{CO}_2$ (requires Pyridoxal Phosphate - Vitamin B6).
    \item $\text{Dopamine} \xrightarrow{\text{Dopamine $\beta$-Hydroxylase}} \text{\textbf{Norepinephrine}}$ (requires Vitamin C and $Cu^{2+}$; occurs inside storage vesicles).
    \item $\text{Norepinephrine} \xrightarrow{\text{PNMT}} \text{\textbf{Epinephrine}}$ (Phenylethanolamine N-methyltransferase uses S-Adenosylmethionine / SAM as methyl donor; stimulated by cortisol in adrenal medulla).
\end{enumerate}
\end{ansbox}

\begin{qbox}{Question 8 [3 $\times$ 3 = 9 Marks]}
Write short notes on:
(a) Cardiac markers in myocardial infarction \\
(b) Hemoglobinopathies: Sickle cell anemia vs. Thalassemia \\
(c) Watson-Crick model of DNA structure
\end{qbox}

\begin{ansbox}
\textbf{(a) Cardiac Markers in Myocardial Infarction (MI):}
\begin{itemize}[leftmargin=*]
    \item Biomarkers released into peripheral circulation following ischemic myocardial cell membrane necrosis:
    \begin{enumerate}
        \item \textbf{Cardiac Troponins (cTnI and cTnT) --- Gold Standard:}
        \begin{itemize}
            \item Highest clinical sensitivity and specificity for myocardial injury.
            \item *Timeline:* Rise within **3--6 hours**, peak at **12--24 hours**, and remain elevated for **7--10 days** (cTnI) or **10--14 days** (cTnT). Useful for delayed diagnosis.
        \end{itemize}
        \item \textbf{Creatine Kinase-MB Isoenzyme (CK-MB):}
        \begin{itemize}
            \item Specific for cardiac tissue ($>6\%$ of total CK).
            \item *Timeline:* Rises in **4--6 hours**, peaks at **18--24 hours**, and normalizes rapidly within **48--72 hours**.
            \item *Unique Clinical Value:* **Detection of early re-infarction** (a secondary rise after 3 days indicates a new infarct).
        \end{itemize}
        \item \textbf{Myoglobin:} Earliest marker to rise (within **1--3 hours**); non-specific (present in skeletal muscle); high negative predictive value.
        \item \textbf{Lactate Dehydrogenase (LDH):} Shows a "flipped pattern" ($\text{LDH}_1 > \text{LDH}_2$) peaking at 3--4 days.
    \end{enumerate}
\end{itemize}

\textbf{(b) Sickle Cell Anemia vs. Thalassemia:}
\begin{center}
\begin{tabularx}{\textwidth}{p{3.2cm}X X}
\toprule
\textbf{Parameter} & \textbf{Sickle Cell Anemia (HbS)} & \textbf{$\beta$-Thalassemia} \\
\midrule
\textbf{Defect Nature} & \textbf{Qualitative} hemoglobinopathy (abnormal structural globin chain). & \textbf{Quantitative} hemoglobinopathy (decreased synthesis of normal globin chains). \\
\textbf{Molecular Mutation} & Point mutation (transversion $A \rightarrow T$) in $\beta$-globin gene: **Glutamic acid $\rightarrow$ Valine at codon 6**. & Variable mutations (nonsense, frameshift, promoter) causing reduced ($\beta^+$) or absent ($\beta^0$) $\beta$-chain synthesis. \\
\textbf{Pathogenesis} & Deoxy-HbS polymerizes into long insoluble fibrous polymers, distorting RBCs into rigid sickle shapes $\rightarrow$ vaso-occlusive crisis. & Excess unpaired $\alpha$-globin chains precipitate, damaging RBC precursors $\rightarrow$ severe ineffective erythropoiesis and hemolysis. \\
\textbf{Red Cell Indices} & Normocytic normochromic anemia with sickled RBCs on smear. & Marked **Microcytic Hypochromic** anemia with target cells and basophilic stippling. \\
\bottomrule
\end{tabularx}
\end{center}

\textbf{(c) Watson-Crick Model of DNA Double Helix (1953):}
\begin{itemize}[leftmargin=*]
    \item \textbf{Double Helical Architecture:} Two polynucleotide chains wound around a central axis to form a right-handed B-DNA double helix.
    \item \textbf{Antiparallel Orientation:} One strand runs in the $5' \rightarrow 3'$ direction; the complementary strand runs in the $3' \rightarrow 5'$ direction.
    \item \textbf{Backbone vs. Core:} Hydrophilic sugar-phosphate backbones form the outside of the helix; hydrophobic nitrogenous bases stack inside perpendicular to axis.
    \item \textbf{Complementary Base Pairing (Chargaff's Rule):}
    \begin{itemize}
        \item Adenine (A) pairs with Thymine (T) via **2 Hydrogen bonds** ($A = T$).
        \item Guanine (G) pairs with Cytosine (C) via **3 Hydrogen bonds** ($G \equiv C$).
    \end{itemize}
    \item \textbf{Dimensions:} Diameter of helix is **2.0 nm** (20 \AA); one complete helical turn spans **3.4 nm** (34 \AA) containing **10 base pairs** (3.4 \AA\ rise per base pair).
    \item Displays alternating **Major Grooves** (2.2 nm) and **Minor Grooves** (1.2 nm) providing access for regulatory DNA-binding proteins.
\end{itemize}
\end{ansbox}

\newpage
\section{Frequency-Based Question Classification (2015--2022)}

\subsection{Very Frequently Repeated Topics ($\ge 80\%$ of Papers)}
\begin{itemize}[leftmargin=*]
    \item \textbf{Oxidative Phosphorylation \& ETC (91\%):} Complexes I--IV, Chemiosmotic theory, proton-motive force, ATP synthase mechanism, uncouplers (DNP), and inhibitors.
    \item \textbf{Lipoprotein Metabolism (82\%):} Classification (Chylomicrons, VLDL, LDL, HDL), apolipoproteins, atherogenesis, reverse cholesterol transport.
    \item \textbf{Enzyme Kinetics \& Inhibition (82\%):} Coenzymes, isoenzymes, $K_m$ and $V_{max}$, competitive vs non-competitive inhibition, Lineweaver-Burk plots.
\end{itemize}

\subsection{Frequently Repeated Topics (60\%--79\% of Papers)}
\begin{itemize}[leftmargin=*]
    \item \textbf{Iron Metabolism (64\%):} Non-heme iron absorption, DMT-1, ferroportin, transferrin, ferritin, hemosiderin, hepcidin regulation.
    \item \textbf{Free Radicals \& Antioxidant Defense (64\%):} ROS generation, superoxide dismutase, catalase, glutathione peroxidase, Vitamin C and E.
    \item \textbf{Protein Structure (64\%):} Primary, secondary ($\alpha$-helix, $\beta$-sheet), tertiary, and quaternary structural bonds and organization.
    \item \textbf{Neurotransmitters \& Catecholamine Biosynthesis (64\%):} Synthesis of DOPA, dopamine, norepinephrine, epinephrine, and degradation.
\end{itemize}

\subsection{Moderately Repeated Topics (40\%--59\% of Papers)}
\begin{itemize}[leftmargin=*]
    \item \textbf{TCA Cycle (55\%):} Steps, enzymes, energetics (10 ATP), amphibolic nature, common terminal pathway.
    \item \textbf{Cardiac Markers (55\%):} Troponin I/T, CK-MB, myoglobin kinetics in myocardial infarction.
    \item \textbf{Heme Biosynthesis \& Porphyrias (55\%):} ALA synthase rate-limiting step, lead poisoning, acute intermittent porphyria.
    \item \textbf{Ketone Body Metabolism (45\%):} HMG-CoA synthase, ketogenesis in liver, ketolysis in extrahepatic tissues, diabetic ketoacidosis.
    \item \textbf{Glycolysis Steps \& Regulation (45\%):} Hexokinase, PFK-1, pyruvate kinase, Rapoport-Luebering cycle in RBCs (2,3-BPG).
    \item \textbf{Hemoglobinopathies (36\%):} Sickle cell anemia molecular mutation vs $\beta$-thalassemia.
\end{itemize}

\subsection{Rarely Repeated Topics (20\%--39\% of Papers)}
\begin{itemize}[leftmargin=*]
    \item \textbf{Fatty Acid Oxidation (30\%):} Carnitine shuttle, $\beta$-oxidation steps (dehydrogenation, hydration, oxidation, thiolysis).
    \item \textbf{Vitamins (30\%):} Vitamin B12 and folate trap, Vitamin C collagen hydroxylation in scurvy.
    \item \textbf{Watson-Crick DNA Double Helix (27\%):} Antiparallel strands, base pairing rules, dimensions.
    \item \textbf{HMP Shunt / Pentose Phosphate Pathway (27\%):} G6PD, NADPH generation, glutathione reduction.
\end{itemize}

\subsection{Unique / One-Time Questions ($<20\%$ of Papers)}
\begin{itemize}[leftmargin=*]
    \item Nucleotide salvage pathway and Lesch-Nyhan syndrome.
    \item Mineral metabolism: Calcium and Phosphorus regulation by PTH, Calcitriol, Calcitonin.
    \item Biochemical basis of Organophosphorus poisoning (Acetylcholinesterase inhibition).
\end{itemize}

\newpage
\appendix
\section{Biochemistry Quick Revision Cheatsheet}

\subsection{Carbohydrate Metabolism Summary}
\begin{itemize}[leftmargin=*]
    \item \textbf{Glycolysis Irreversible Steps:} 1. Glucose $\rightarrow$ G-6-P (Hexokinase/Glucokinase); 2. F-6-P $\rightarrow$ F-1,6-bisP (PFK-1, **Rate-Limiting**); 3. PEP $\rightarrow$ Pyruvate (Pyruvate kinase).
    \item \textbf{Rapoport-Luebering Cycle (RBCs):} Bypasses PGK to make **2,3-BPG**; shifts $\text{O}_2$-dissociation curve to the RIGHT (promotes $\text{O}_2$ release to tissues).
    \item \textbf{HMP Shunt:} Yields **NADPH** (for biosynthetic reactions and maintaining reduced glutathione against oxidative hemolysis) and Ribose-5-P (for nucleotide synthesis).
    \item \textbf{TCA Cycle:} Yields $3\ \text{NADH}, 1\ \text{FADH}_2, 1\ \text{GTP} \rightarrow 10\ \text{ATP}$ per Acetyl-CoA. Amphibolic: provides precursors for heme (Succinyl-CoA), glucose/aspartate (OAA), and fatty acids (Citrate).
\end{itemize}

\subsection{Lipids \& Bioenergetics Summary}
\begin{itemize}[leftmargin=*]
    \item \textbf{Ketone Bodies:} Acetoacetate, $\beta$-hydroxybutyrate, Acetone. Liver synthesizes them (HMG-CoA synthase) but **cannot utilize them** (lacks thiophorase). Increased in starvation and DKA.
    \item \textbf{Lipoproteins:} Chylomicrons (dietary TG); VLDL (hepatic endogenous TG); LDL (cholesterol to tissues, atherogenic); HDL (reverse cholesterol transport, cardioprotective).
    \item \textbf{ETC Complexes:} I (NADH-Q, pumps $4\ \text{H}^+$), II (Succinate-Q, no $\text{H}^+$ pump), III (Cytochrome $bc_1$, pumps $4\ \text{H}^+$), IV (Cytochrome oxidase, pumps $2\ \text{H}^+$).
    \item \textbf{Chemiosmotic Theory:} Proton-motive force drives ATP synthesis via Complex V ($F_0F_1$ ATP synthase). Uncouplers (2,4-DNP, thermogenin) dissipate gradient as heat.
\end{itemize}

\subsection{Proteins, Enzymes \& Clinical Chemistry}
\begin{itemize}[leftmargin=*]
    \item \textbf{Enzyme Inhibition:}
    \begin{itemize}
        \item *Competitive:* Inhibitor binds active site; **$V_{max}$ unchanged, $K_m$ increased**. Overcome by high $[S]$ (e.g. Statins).
        \item *Non-Competitive:* Inhibitor binds allosteric site; **$V_{max}$ decreased, $K_m$ unchanged** (e.g. Cyanide).
    \end{itemize}
    \item \textbf{Iron Homeostasis:} Absorbed as $\text{Fe}^{2+}$ via DMT-1; transported as $\text{Fe}^{3+}$ on Transferrin; stored in Ferritin. **Hepcidin** degrades ferroportin to block iron export.
    \item \textbf{Cardiac Markers:} **cTnI/cTnT** (rises 3--6h, peaks 24h, stays elevated 10--14d); **CK-MB** (rises 4--6h, normalizes in 48--72h; ideal for re-infarction).
    \item \textbf{Sickle Cell Mutation:} $\beta$-globin codon 6: Glutamic acid $\rightarrow$ Valine.
\end{itemize}

\vspace{1cm}
\begin{center}
\textbf{--- End of Biochemistry Solutions ---}
\end{center}

\end{document}
